%%
%% Energy-aware scheduling in the interval job shop:
%% from lexicographic to Pareto optimisation of makespan and idle energy
%%
%% Working draft — rev. 2026-07-13. Built on Elsevier's elsarticle template
%% (same conventions as the ASOC tabu manuscript).
%% Compile with:  latexmk -pdf main
%%
%% STATUS: all four experiments complete on the full 82-instance benchmark
%% (lexicographic, right-shift, N2ME, capped ladder vs. dominance controls).
%% Author, funding and figures are in place. Remaining \todo{...} boxes are
%% all OPTIONAL/pending-a-decision: a Zenodo data deposit (DOI), and two
%% supplementary items (epsilon-indicator vs. exact fronts, tuning-gain table)
%% plus optional strengthening of the conflict probe and validation fronts.
%%
\documentclass[preprint,12pt]{elsarticle}

\usepackage{amssymb}
\usepackage{amsmath}
\usepackage{amsthm}
\usepackage{xcolor}
\usepackage{hyperref}
\usepackage{booktabs}
\usepackage{algorithm}
\usepackage{algpseudocode}
\usepackage{float}
\usepackage{pgfplots}
\pgfplotsset{compat=1.15}

\newtheorem{theorem}{Theorem}
\newtheorem{proposition}{Proposition}
\newtheorem{corollary}{Corollary}
\newtheorem{definition}{Definition}
\newtheorem{lemma}{Lemma}
\newtheorem{remark}{Remark}

\newcommand{\todo}[1]{\noindent\fcolorbox{red}{yellow!20}{\parbox{0.97\linewidth}{\textbf{TODO:} #1}}}
\newcommand{\Cmax}{C_{\max}}
\newcommand{\NPE}{\mathrm{NPE}}
\newcommand{\lexle}{\preceq_{\mathrm{LEX2}}}

\begin{document}

\begin{frontmatter}

\title{Energy-aware scheduling in the interval job shop:\\
from lexicographic to Pareto optimisation of makespan and idle energy}

\author[uniovi]{Hern\'an D\'iaz Rodr\'iguez\corref{cor1}}
\ead{diazhernan@uniovi.es}
\cortext[cor1]{Corresponding author}
\address[uniovi]{Department of Computer Science, University of Oviedo, Campus
of Gij\'on, 33204 Gij\'on, Spain}

\begin{abstract}
We study, for the first time, the multiobjective interval job shop problem
minimising interval makespan and non-processing (idle) energy. In the
non-flexible job shop, active energy is schedule-invariant, so the green
objective reduces exactly to idle passive energy; moreover, the apparent
makespan--energy correlation is an artefact of semi-active timing --- idle
energy is not regular, and exact $\varepsilon$-constraint fronts reveal a
steep conflict once start times may shift. The decision space is therefore
the pair (order, timing). On a benchmark of 82 instances extended with power
data, we develop: a lexicographic tie-break that harvests the energy hidden
in makespan plateaus (7.9\% mean idle-energy reduction, up to 26.5\%, at no
makespan cost); an exact and a near-optimal heuristic right-shift timing
post-optimisation (a further ${\sim}8\%$, complementary); a novel
bi-critical neighbourhood targeting idle gaps, with feasibility and
connectivity results; and, for the Pareto front, a capped
$\varepsilon$-ladder that recycles the single-objective tabu solver,
unchanged, into a front generator. Against irace-tuned NSGA-II and
Pareto-archive controls, the ladder attains the highest hypervolume on 55 of
82 instances and dominates up to size $20{\times}20$; we delineate --- and
explain mechanistically --- the large-instance regime where the scalar
anchor loses traction and dominance-based search becomes competitive again.
\end{abstract}

\begin{keyword}
Job shop scheduling \sep Interval uncertainty \sep Green scheduling \sep
Multiobjective optimisation \sep $\varepsilon$-constraint \sep
Pareto front \sep Local search
\end{keyword}

\end{frontmatter}

%=============================================================================
\section{Introduction}
\label{sec:intro}

Production scheduling has traditionally optimised time-based performance
measures, most prominently the makespan. Two independent concerns have
reshaped the field in the last decade. First, the environmental impact and
cost of energy have promoted \emph{green scheduling}, where energy
consumption co-exists with classical objectives
\cite{Afsar2022,GonzalezRodriguez2020}. Second, the acknowledgement that
activity durations are rarely known in advance has motivated models of
uncertainty; intervals are arguably the least committing such model,
requiring only a lower and an upper bound for each duration
\cite{Diaz2020,Diaz2023FEABC}.

Energy-aware machine scheduling is by now a mature field: the survey of
\cite{Gahm2016} organises it by energy lever (idle reduction, speed
scaling, on/off policies) and observes that idle energy is the lever most
amenable to pure sequencing decisions. For the deterministic job shop,
multiobjective genetic algorithms trading production objectives against
energy are well established \cite{May2015,ZhangChiong2016}. Under
uncertainty, both concerns have been combined for the \emph{flexible} job
shop with interval durations, where a lexicographic goal-programming
strategy minimises makespan first and total energy second
\cite{Afsar2024,Afsar2025}, and for the \emph{fuzzy} job shop, where
Pareto-based memetic algorithms optimise makespan and non-processing
energy simultaneously \cite{Afsar2022,GonzalezRodriguez2020}. However, to
the best of our knowledge, the multiobjective \emph{interval} job shop ---
without flexibility --- has not been addressed: the energy-aware
neighbourhood of \cite{GarciaGomez2023} relies on machine reassignment
moves that require flexibility, and its passive-energy component reduces
energy only through the makespan. This paper fills that gap, and argues
that the non-flexible interval case is not a mere restriction: it has
structure of its own that changes both the model and the search operators.

On the methodological side, the two classical routes to a Pareto front are
dominance-based evolutionary search, epitomised by NSGA-II \cite{Deb2002},
and scalarising decompositions --- the $\varepsilon$-constraint method
\cite{Haimes1971,Miettinen1999} and, in the evolutionary setting, MOEA/D
\cite{ZhangLi2007}. Our front generator belongs squarely to the second
family; front quality is assessed with the hypervolume and the additive
$\varepsilon$-indicator, following the recommendations of
\cite{Zitzler2003}.

Our contributions are the following:
\begin{enumerate}
\item \textbf{Model.} We show that in the non-flexible job shop the active
energy is schedule-invariant (Proposition~\ref{prop:ae}), so the green
objective reduces exactly to the passive energy of idle intervals, for which
we give an extreme-scenario definition evaluated component-wise
(Section~\ref{sec:problem}).
\item \textbf{Conflict analysis.} We show that under semi-active timing
makespan and idle energy are strongly correlated (Pearson $+0.80$ to
$+0.96$), which would render the biobjective problem nearly trivial --- and
that this is an artefact: idle energy is not a regular measure, and exact
$\varepsilon$-constraint fronts obtained with constraint programming exhibit
a genuine trade-off once start times may shift (idle energy drops by
${\sim}56\%$ when the makespan is relaxed by 10\%). The decision space of
the green JSP is therefore the pair \emph{(order, timing)}
(Section~\ref{sec:conflict}).
\item \textbf{Benchmark.} We extend the 82-instance IJSP benchmark with
machine power data under a documented, reproducible scheme
(Section~\ref{sec:benchmark}).
\item \textbf{Methods.} We propose a family of methods on top of a
memetic algorithm with tabu search local search
(Section~\ref{sec:methods}): for a single compromise solution, a
lexicographic makespan-then-energy strategy, an exact and a heuristic
right-shift timing post-optimisation, and a novel \emph{bi-critical
neighbourhood} $N_2^{ME}$ combining makespan-critical and energy-critical
arcs, with feasibility proven and bi-objective connectivity analysed; and,
for the whole front, a \emph{capped $\varepsilon$-ladder} that reuses the
strong single-objective interval-JSP tabu search --- unchanged --- as a
Pareto front generator, by decomposing the front into warm-started,
makespan-capped energy problems. The decomposition itself is not new: the
ladder is the classical $\varepsilon$-constraint method~\cite{Haimes1971},
in the decomposition spirit of MOEA/D~\cite{ZhangLi2007}, specialised to this
problem. What is ours is the instantiation --- recycling the tuned solver
unchanged through a clamped-LEX2 fitness --- and the structural guarantees
it carries for the idealised (exact-per-level) ladder: goal-programming and
$\varepsilon$-constraint equivalence, plus capped connectivity of the
neighbourhood; the metaheuristic realisation inherits the structure, not an
approximation bound.
\item \textbf{When decomposition beats dominance, and why.} Rather than
crown a ``best'' multiobjective metaheuristic, we use two off-the-shelf
dominance-based algorithms (an NSGA-II memetic and a Pareto-archive ABC) as
\emph{controls} to test whether the front comes for free, and we characterise
the outcome: on the 82-instance benchmark decomposition dominates on small
and medium instances and in aggregate, while a landscape analysis explains
both why (energy-deep regions lie behind makespan valleys that dominance
cannot cross) and where it breaks down (large high-ratio instances, where the
scalar anchor loses traction and dominance recovers)
(Section~\ref{sec:experiments}).
\end{enumerate}

Taken together, the contribution is not another multiobjective solver that
edges out its predecessors: it is a new problem, a method that recycles an
existing single-objective solver into a Pareto front generator with
guarantees, and a mechanism-level account of when that strategy is the right
one. The dominance-based algorithms are the yardstick that makes the last two
claims credible, not the object of a horse race.

%=============================================================================
\section{The interval job shop with idle energy}
\label{sec:problem}

\subsection{Interval JSP}

An instance consists of $n$ jobs, each a chain of $m$ tasks, every task
requiring a specific machine among $m$ machines; each machine processes one
task at a time and no two tasks of a job overlap. The processing time of
task $x$ is only known to lie in an interval
$\mathbf{p}_x=[p^-_x,p^+_x]$. A solution is a processing order $\sigma$ (a
machine sequencing compatible with job precedences), which orients the
disjunctive arcs $D(\sigma)$ of the solution graph $G(\sigma)$.

Following \cite{DiazTabu2026}, the two \emph{extreme graphs} $G^-(\sigma)$
and $G^+(\sigma)$ weight every arc $(x,y)$ with $p^-_x$ and $p^+_x$
respectively. For $i\in\{-,+\}$ the \emph{head} $r^i_v$ of task $v$ is its
earliest starting time in $G^i(\sigma)$ (\emph{semi-active timing}),
computed component-wise:
\begin{equation}
r^i_v=\max\bigl\{\,r^i_{JP(v)}+p^i_{JP(v)},\;
r^i_{MP_\sigma(v)}+p^i_{MP_\sigma(v)}\,\bigr\},
\label{eq:heads}
\end{equation}
where $JP(v)$ and $MP_\sigma(v)$ are the job and machine predecessors of
$v$ (missing predecessors contribute $0$). The interval makespan is
$\mathbf{C}_{\max}(\sigma)=[C^-_{\max},C^+_{\max}]$ with
$C^i_{\max}=\max_v (r^i_v+p^i_v)$.

Interval objectives require a ranking. We adopt LEX2, the lexicographic
order that compares upper bounds first and breaks ties with lower bounds;
in robust-optimisation terms LEX2 performs lexicographic min--max
optimisation, and it was the best-performing ranking in the extensive
comparison of \cite{DiazTabu2026}.

\subsection{Energy model}

We adopt the energy model of the interval green FJSP literature
\cite{Afsar2024,Afsar2025}: machine $k$ is switched on when its first task
starts and off when its last task completes; while on, it draws a
\emph{passive power} $P_k$ per time unit, plus an \emph{active} power while
actually processing. The total energy is
$TE=\sum_k (PE_k+AE_k)$ with $AE_k=\sum_{x\in k} Pa_{x}\,\mathbf{p}_x$.

\begin{proposition}[Active energy invariance]
\label{prop:ae}
In the non-flexible job shop, $AE=\sum_k AE_k$ does not depend on the
processing order $\sigma$.
\end{proposition}
\begin{proof}
Without flexibility every task is always processed on the same machine with
the same (interval) duration, so each term $Pa_x\,\mathbf{p}_x$ is
instance data. \end{proof}

Consequently the only energy term the scheduler controls is the passive
energy of \emph{idle} intervals, the non-processing energy. For machine $k$
let $u^k_1\prec\cdots\prec u^k_{n_k}$ be its processing sequence under
$\sigma$. The component-$i$ idle gap between consecutive tasks is
$\gamma^i(u^k_j,u^k_{j+1})=r^i_{u^k_{j+1}}-(r^i_{u^k_j}+p^i_{u^k_j})\ge 0$
and
\begin{equation}
\NPE^i(\sigma)=\sum_{k} P_k \sum_{j<n_k}\gamma^i(u^k_j,u^k_{j+1})
=\sum_{k} P_k\bigl(C^i(k)-S^i(k)\bigr)-\beta^i,
\label{eq:npe}
\end{equation}
where $C^i(k)$ and $S^i(k)$ are the machine's last completion and first
start and $\beta^i=\sum_k P_k\sum_{x\in k}p^i_x$ is a constant. The
\emph{gap form} = \emph{window form} identity \eqref{eq:npe} is used
throughout: NPE improves either by finishing a machine earlier or by
switching it on later. The energy objective is the pair of extreme-scenario
evaluations $\mathbf{\NPE}(\sigma)=(\NPE^-,\NPE^+)$, ranked by LEX2 as well
(upper scenario first, lower on ties).

\begin{remark}[$\mathbf{\NPE}$ is a scenario pair, not an ordered interval]
\label{rem:npe-pair}
Unlike the makespan, whose components satisfy $C^-_{\max}\le C^+_{\max}$
because the longest path is monotone in the durations, the idle gaps are
\emph{not} monotone, and $\NPE^->\NPE^+$ is possible. A two-machine example:
let machine $k$ process $a$ then $b$, where $a$ starts at $0$ with
$\mathbf{p}_a=[2,9]$ and $b$ waits for its job predecessor $c$ (on the other
machine, starting at $0$) with crisp $\mathbf{p}_c=[10,10]$. The gap before
$b$ is $\max(0,\,p^i_c-p^i_a)$: $8$ in the lower scenario but $1$ in the
upper. We therefore write $\mathbf{\NPE}$ as a pair $(\NPE^-,\NPE^+)$ ---
each component is the exact idle energy of one extreme scenario --- and rank
pairs lexicographically by LEX2, which retains its min--max reading
(optimise the worst-case scenario first). None of the results in this paper
require $\NPE^-\le\NPE^+$; dominance and the archive operate on the pair
order directly.
\end{remark}

\begin{remark}[Non-regularity]
\label{rem:nonregular}
$\NPE$ is not a regular performance measure \cite{Afsar2022}: delaying a
task can reduce it. Unless stated otherwise we evaluate $\NPE$ on the
semi-active schedule \eqref{eq:heads}, which makes it a function of
$\sigma$ alone; Section~\ref{sec:rs} restores the timing dimension.
\end{remark}

%=============================================================================
\section{Is there a trade-off? Conflict analysis}
\label{sec:conflict}

A biobjective problem is only meaningful if its objectives conflict
\cite{Afsar2022}. A first experiment suggests they do not: evaluating
makespan and semi-active NPE on $200$ random schedules per instance yields
Pearson correlations of $+0.80$ (\texttt{ft10}), $+0.94$ (\texttt{la29})
and $+0.96$ (\texttt{tai30\_20\_01}); on near-optimal schedules the
apparent Pareto front collapses to $1$--$4$ points. Compressing the
makespan compresses machine windows, and semi-active NPE follows suit.

This correlation is an artefact of semi-active timing. Because NPE is not
regular (Remark~\ref{rem:nonregular}), the proper decision space is the
pair (order, timing). We computed exact fronts by
$\varepsilon$-constraint with CP-SAT on midpoint durations, with free
start times: minimise NPE subject to $\Cmax\le(1+\varepsilon)\,C^*$,
sweeping $\varepsilon$ from $0$ to $10\%$.
Table~\ref{tab:probe} shows a genuine and steep trade-off: a 10\% makespan
allowance halves the idle energy.

\begin{table}[htb]
\centering\small
\caption{Exact $\varepsilon$-constraint probe (CP-SAT, midpoint durations,
free start times, 60\,s per point). NPE at selected makespan allowances.}
\label{tab:probe}
\begin{tabular}{l rrrr r}
\toprule
Instance & $\varepsilon{=}0$ & $+3\%$ & $+5\%$ & $+10\%$ & drop \\
\midrule
ft10 (10$\times$10) & 5431 & 3305 & 2819 & 2400 & $-55.8\%$ \\
la29 (20$\times$10) & --\,\footnotemark[1] & --\, & 3230 & 2480 & $-55.3\%$ \\
\bottomrule
\end{tabular}

\footnotesize $^1$ CP-SAT could not close the tightest allowances on
20$\times$10 within the budget; the drop is measured from the tightest
solved point.
\end{table}

Two design consequences follow. First, algorithms must include a timing
component (Section~\ref{sec:rs}); a search that only manipulates the order
under semi-active evaluation can at most exploit the makespan plateaus.
Second, exact fronts do not scale beyond small instances, so reference
fronts for larger instances must be approximated by the algorithms
themselves, as in \cite{Afsar2022}.

\todo{Extend Table~\ref{tab:probe} with 2--3 additional small instances and
a finer $\varepsilon$ grid (probe script ready:
\texttt{experiments/mo\_green\_2026/epsilon\_probe.py}).}

%=============================================================================
\section{Benchmark with energy data}
\label{sec:benchmark}

We use the 82 interval instances of \cite{DiazTabu2026}: 12 classical
OR-Library instances and 70 Taillard-derived interval instances
(\texttt{tai15\_15} to \texttt{tai50\_20}). Since no power data exists for
them, we extend every instance with integer passive powers
$P_k\sim U\{2,\dots,8\}$, drawn per instance with a fixed documented seed,
mirroring the way the fuzzy FJSP benchmark was extended with energy data
\cite{GarciaGomez2023}. Active powers are not needed by
Proposition~\ref{prop:ae}. The extended instances and the generator script
are publicly available.\footnote{Zenodo DOI to be minted with the final
dataset; generator \texttt{extend\_instances.py}, seed $23$.}

%=============================================================================
\section{Solution methods}
\label{sec:methods}

All methods build on the memetic algorithm of \cite{DiazTabu2026}: an
artificial-bee-colony-inspired evolutionary algorithm (ABC-PSO) hybridised
with tabu search local search over the makespan neighbourhood $N_2$
(reversal of boundary arcs of extreme critical blocks), with irace-tuned
hyperparameters. We keep those hyperparameters untouched in all arms, so
every observed difference is attributable to the multiobjective machinery.

\subsection{Lexicographic makespan-then-energy}
\label{sec:lex}

The first rung replaces the fitness by the lexicographic pair
$\bigl(\mathbf{C}_{\max},\mathbf{\NPE}\bigr)$: candidates are compared by
interval makespan under LEX2 and, only on ties, by the NPE scenario pair
under LEX2
--- a four-level comparison
$(C^+_{\max},C^-_{\max},\NPE^+,\NPE^-)$. The comparison is applied in
\emph{every} decision of the algorithm (selection, replacement, elitism,
local search acceptance), so the pressure towards low energy acts
throughout the search, but only inside makespan plateaus. The tabu search
itself remains single-objective on the primary component, following the
design principle of one intensifier per objective of \cite{Afsar2022}; the
full lexicographic fitness of the improved solution is rebuilt after the
local search returns.

The rationale is that JSP makespan plateaus are huge: a makespan-only
search drifts through them blindly, ending anywhere; the lexicographic
tie-break walks them towards their lowest-energy point at zero cost in the
primary objective. Section~\ref{sec:exp-lex} quantifies exactly how much
energy those plateaus contain.

\subsection{Right-shift timing optimisation}
\label{sec:rs}

The second rung restores the timing dimension. For a \emph{fixed} order
$\sigma$ and fixed per-component makespans, the optimal timing problem is:
delay task start times, respecting all precedences and without exceeding
$C^i_{\max}(\sigma)$, so as to minimise \eqref{eq:npe}. By the window
identity the objective is linear in the start times and the constraints
form a difference system, so the problem is a linear programme with
integral optima, solved exactly in milliseconds per solution; this is the
per-order analogue of the LP post-processing of \cite{Afsar2022}. We also
define a single-pass heuristic in the spirit of their right-shift
procedure: machine-last tasks are pinned to their semi-active starts
(delaying them can only widen the window), and every other task, in
reverse topological order, is delayed up to the minimum start of its
successors. The heuristic is $O(nm)$ and captured 86--95\% of the exact
gain in our validation. Both variants preserve
$\mathbf{C}_{\max}$ by construction and are applied independently per
interval component.

\subsection{The bi-critical neighbourhood $N_2^{ME}$}
\label{sec:n2me}

The third rung makes the \emph{search} energy-aware. Existing energy
neighbourhoods rely on flexibility (reassignment moves) or reduce passive
energy only through the makespan \cite{GarciaGomez2023}; to the best of our
knowledge, no sequencing neighbourhood targets idle gaps. We propose one,
built on the $H(\sigma)$ framework of \cite{DiazTabu2026}.

\begin{definition}[Tight arc, sustaining chain]
An arc $(x,y)$ of $G(\sigma)$ is \emph{$i$-tight} if $r^i_y=r^i_x+p^i_x$.
A \emph{sustaining chain} $\pi^i(v)$ is a chain of $i$-tight arcs from a
task with null head to $v$; its length is $r^i_v$. Critical paths are the
sustaining chains of the tasks realising $C^i_{\max}$.
\end{definition}

\begin{definition}[Energy-critical arcs; bi-critical neighbourhood]
\label{def:n2me}
For $i\in\{-,+\}$ let $\mathcal{V}^i(\sigma)$ contain the right endpoints
of positive gaps and the last task of every machine. The energy-critical
arcs $E(\sigma)$ are the disjunctive arcs on sustaining chains of tasks in
$\mathcal{V}^i(\sigma)$. Then
$N_2^{ME}(\sigma)$ reverses the arcs of $N_2(\sigma)\cup E(\sigma)$. In
\emph{lexicographic mode} energy moves are accepted only if the makespan is
not LEX2-worsened; in \emph{Pareto mode} all moves are offered to the
acceptance criterion.
\end{definition}

\begin{theorem}[Feasibility]
The reversal of any $i$-tight disjunctive arc yields a feasible order;
hence every neighbour in $N_2^{ME}(\sigma)$ is feasible.
\end{theorem}
\begin{proof}[Proof idea]
The feasibility proof of \cite{DiazTabu2026} for critical arcs uses only
tightness of the reversed arc; all arcs of $E(\sigma)$ are tight by
definition. \end{proof}

\begin{theorem}[Connectivity]
\label{thm:connectivity}
From every order $\sigma$ there is a finite sequence of $N_2^{ME}$ moves
reaching an order that is optimal in makespan and, among those, minimal in
every machine completion $C^i(k)$.
\end{theorem}
\begin{proof}[Proof idea]
Makespan deficits are repaired by the connectivity theorem of
\cite{DiazTabu2026} ($N_2\subseteq N_2^{ME}$). For a completion deficit at
equal makespan, a sustaining chain of the machine-last task with all arcs
agreeing with the target would certify an equal completion in the target
--- a contradiction; hence some disagreeing arc lies on the chain, i.e.\ in
$E(\sigma)$, and the inversion-count argument of \cite{DiazTabu2026}
applies. \end{proof}

\begin{remark}
Whether connectivity extends to the machine \emph{start} terms of
\eqref{eq:npe} (switching machines on later) is open; candidate delay
moves and the formal statement are discussed in the companion technical
note. In practice the right-shift operator of Section~\ref{sec:rs}
optimises exactly those terms for each visited order, so the search and
the timing operator are complementary by design.
\end{remark}

\begin{remark}[Scale behaviour]
\label{rem:n2me-scale}
The experimental ablation (Section~\ref{sec:exp-n2me}) shows that making
the tabu search energy-aware pays off only up to medium-sized instances:
on large instances the energy moves crowd out the makespan-worsening
moves the tabu search needs for diversification, and quality degrades in
both objectives. Candidate-set restriction and lazy energy evaluation
mitigate but do not remove the effect; gating the energy arcs to
incumbent-makespan plateaus is left as an open design.
\end{remark}

\subsection{Front generation: dominance-based baselines as controls}
\label{sec:pareto-baselines}

The natural, off-the-shelf way to return a front is a dominance-based
multiobjective evolutionary algorithm, and this is exactly what a
practitioner would reach for. We therefore use two such algorithms not as
competitors to be defeated but as \emph{controls}: they measure how much of
the front a standard dominance-based search recovers on this problem, and
hence whether the proposed decomposition is needed at all. To keep the
comparison about the front-generation \emph{strategy} and nothing else, both
baselines share the makespan-energy evaluation, the $N_2^{ME}$ local search,
the time budget, and the per-arm irace tuning of the proposed method; only
the way the front is assembled differs. Dominance is defined over the product
of the two LEX2 total orders: $\sigma$ dominates $\sigma'$ iff
$\mathbf{C}_{\max}(\sigma)\lexle\mathbf{C}_{\max}(\sigma')$ and
$\mathbf{\NPE}(\sigma)\lexle\mathbf{\NPE}(\sigma')$, one of them strictly
(the interval-objective adaptation of standard dominance). We instantiate
two such baselines, both sharing the makespan-energy evaluation and the
$N_2^{ME}$ local search: \textbf{ABC-P}, the ABC-PSO engine of
\cite{DiazTabu2026} with a bounded non-dominated archive fed each
generation; and \textbf{MA-P}, a memetic engine with true NSGA-II survival
\cite{Deb2002} (non-dominated sorting plus crowding distance in the
replacement), following the fuzzy-JSP design of \cite{Afsar2022}. The archive stores the
non-dominated solutions found; capacity is controlled by crowding distance.

\subsection{Front generation: the capped $\varepsilon$-ladder}
\label{sec:pareto}

Our proposed method exploits a structural observation, made precise by the
one-move landscape analysis of Section~\ref{sec:exp-exact}: from a
makespan-optimal anchor, the energy-deep region of the front is not
reachable by single-move local search --- it lies behind makespan valleys
that only recombination crosses. Rather than ask one population to cover
the whole front at once, the ladder decomposes it into single-objective
levels and concentrates the search on one level at a time.

Methodologically this is the classical $\varepsilon$-constraint
method~\cite{Haimes1971}, and shares the decompose-into-scalar-subproblems
philosophy of MOEA/D~\cite{ZhangLi2007} and of the lexicographic
goal-programming used for the interval green FJSP~\cite{Afsar2024,Afsar2025};
we make no claim to the decomposition idea itself. Our specific contribution
is a realisation tailored to the interval JSP: the makespan bound is imposed
not as a hard constraint but as a \emph{clamped} objective (below), so that
the tuned single-objective tabu search --- operators, pruning, and $N_2^{ME}$
neighbourhood alike --- is reused verbatim; the interval nature of the
objectives introduces a soundness subtlety absent from the scalar case, which
we resolve at the archive (Remark~\ref{rem:soundness}); and the neighbourhood
retains a connectivity guarantee under the cap
(Theorem~\ref{thm:capped-conn}).

\begin{algorithm}[H]
\caption{Capped $\varepsilon$-ladder (per run).}
\begin{algorithmic}[1]
\State \textbf{Anchor:} run the (random-initialised) engine minimising
$\mathbf{C}_{\max}$ lexicographically; let $C^{*}$ be its makespan and seed
the archive with its solutions.
\For{$\varepsilon$ in an increasing grid}
  \State $\kappa \gets (1+\varepsilon)\,C^{*}$ \Comment{makespan cap}
  \State run the engine ranking by the \emph{clamped pair}
  $\bigl(\max(\mathbf{C}_{\max},\kappa),\mathbf{\NPE}\bigr)$, warm-started
  from the previous level's front $+$ the anchor
  \State add the level's non-dominated solutions (unclamped) to the archive
\EndFor
\State right-shift (Section~\ref{sec:rs}) every archived solution; return
the archive
\end{algorithmic}
\end{algorithm}

Clamping the makespan to the cap makes all cap-feasible solutions tie on
the primary objective, so the entire evolutionary pressure --- selection,
survival, and the $N_2^{ME}$ tabu search alike --- is redirected onto
energy, without changing any operator (Proposition~\ref{prop:goal}: each
level is a lexicographic goal-programming problem with the cap as target).
Warm seeding is essential rather than cosmetic: on medium and large
instances a from-scratch level spends its whole budget merely re-reaching
the cap region. The per-run budget is size-adaptive but always within the
same ${\sim}900$\,s total granted to each baseline run: anchor
$400/500/600$\,s plus $8/6/5$ levels of $60$--$65$\,s on
classical-and-small/$30{\times}$/$50{\times}$ instances respectively, so no
arm receives more compute. Section~\ref{sec:ladder} gives the theoretical
backing:
level optima are exactly the Pareto points at their cap
(Proposition~\ref{prop:eps}), and $N_2^{ME}$ retains connectivity under a
cap for the makespan and machine-completion terms (Theorem~\ref{thm:capped-conn}).

Quality is measured by hypervolume ratio and the additive
$\varepsilon$-indicator against a per-instance reference front (the
non-dominated union of all arms), and against the exact fronts of
Section~\ref{sec:conflict} on small instances.

\subsection{Theoretical guarantees of the ladder}
\label{sec:ladder}

Write $\kappa=(1+\varepsilon)C^{*}$ for a level's makespan cap, applied
component-wise to the interval, $\kappa=(\kappa^-,\kappa^+)$, and
$\Sigma_\kappa=\{\sigma: C^-_{\max}(\sigma)\le\kappa^-\wedge
C^+_{\max}(\sigma)\le\kappa^+\}$ for its cap-feasible set. A level ranks
solutions by the clamped pair
$\Phi_\kappa(\sigma)=\bigl(\max(\mathbf{C}_{\max},\kappa),\mathbf{\NPE}\bigr)$,
compared with LEX2 inside each component. Three results, whose full proofs are
given in the supplementary technical note and sketched here, justify the
construction.

\begin{proposition}[Goal-programming form]\label{prop:goal}
Ordering by $\Phi_\kappa$ coincides with the lexicographic goal-programming
order with target $\kappa$: makespan deviation
$D_\kappa=[\max(0,C^--\kappa^-),\max(0,C^+-\kappa^+)]$ first, energy second.
\end{proposition}
\noindent Component-wise $\max(C,\kappa)=\kappa+\max(0,C-\kappa)$, and adding
the constant $\kappa$ preserves every LEX2 comparison; each level is thus the
interval green-FJSP goal program with the cap as the production goal, and the
ladder is that methodology turned into a front generator by sweeping the goal.

\begin{proposition}[$\varepsilon$-constraint soundness and completeness]
\label{prop:eps}
With the level order refined to the lexicographic triple
$(D_\kappa,\mathbf{\NPE},\mathbf{C}_{\max})$ and $\Sigma_\kappa\neq\emptyset$:
(i) every level optimum lies in $\Sigma_\kappa$, LEX2-minimises $\mathbf{\NPE}$
over it, and is non-dominated within $\Sigma_\kappa$; and (ii) every
Pareto-optimal order is the level optimum for the cap equal to its own
makespan, so the ideal ladder over a cap grid $\Gamma$ returns exactly the
Pareto front restricted to $\Gamma$, and grid refinement recovers the whole
front.
\end{proposition}

\begin{remark}[Interval soundness, restored at the archive]
\label{rem:soundness}
Non-dominance ``within $\Sigma_\kappa$'' can fail globally through an interval
subtlety --- a solution whose upper makespan is within the cap but whose lower
is not can be LEX2-better yet clamped-worse. The archive neutralises this by
storing \emph{unclamped} objective values and applying the dominance filter
across all levels, so every reported front point is globally non-dominated;
Proposition~\ref{prop:eps}(ii) then guarantees that no Pareto point is out of
the ladder's reach.
\end{remark}

\begin{theorem}[Connectivity under a cap]\label{thm:capped-conn}
Fix a cap $\kappa$ with $\Sigma_\kappa\neq\emptyset$ and a target
$\sigma^{\circ}\in\Sigma_\kappa$. From any order, finitely many $N_2^{ME}$
moves reach $\Sigma_\kappa$; and from any cap-feasible order, the
machine-completion energy deficits with respect to $\sigma^{\circ}$ are
repairable within $N_2^{ME}$. Machine-start deficits are the one open case,
addressed heuristically by goal-gated delay arcs.
\end{theorem}
\noindent While $\sigma\notin\Sigma_\kappa$, some component has
$C^i_{\max}(\sigma)>\kappa^i\ge C^i_{\max}(\sigma^{\circ})$, so the
makespan-connectivity key lemma of \cite{DiazTabu2026} applies verbatim with
$\sigma^{\circ}$ as target; the completion side is covered unchanged, as it
compares completions rather than makespans. A landscape argument delimits what
connectivity can achieve: an exhaustive one-move probe from a converged
\texttt{ft10} anchor ($89$ reversals, $561$ insertions) found at most two
energy-improving moves against a ${\sim}60\%$ exact-front potential, so
crossing into energy-deep \emph{basins} is the work of recombination in the
population engine, while Theorem~\ref{thm:capped-conn} governs refinement
within the basin already reached. (The probe is a single-instance
motivation, not a general claim; the ladder's design does not depend on
it.)

%=============================================================================
\section{Experimental study}
\label{sec:experiments}

\subsection{Setup}
\label{sec:exp-setup}

All arms run on the 82-instance benchmark with the irace-tuned
configuration of \cite{DiazTabu2026}: 30 runs per instance, 900\,s time
limit, population 250, tabu search on $N_2$. The makespan-only baseline is
\emph{not} re-run: we reuse the published raw results of
\cite{DiazTabu2026} and recompute their NPE uniformly. Both experiments'
solutions are decoded and evaluated by the same independent implementation
(and the same power data), so the energy comparison is
implementation-independent; the limitation is that the baseline's raw
solutions come from a different experimental campaign, so small makespan
differences between campaigns cannot be attributed to the algorithms ---
we therefore never claim a makespan improvement, only the absence of
evidence of a makespan cost.

For Experiment~1 (one value per instance and arm), comparisons use
Wilcoxon signed-rank tests paired by instance. For the three-arm front
comparison of Section~\ref{sec:exp-pareto}, we use a Friedman omnibus test
over the 82 instances followed by pairwise signed-rank post-hoc tests
(paired by instance) with Holm correction; per-instance comparisons of the
30-run hypervolume distributions use unpaired Mann--Whitney $U$ tests
--- runs of different engines are independent, so a per-run pairing would
not be defensible --- Holm-corrected across the 82 instances.

\subsection{How much energy do makespan plateaus contain?}
\label{sec:exp-lex}

Table~\ref{tab:lex} compares the lexicographic arm (LexME) against the
published makespan-only tabu search over the full benchmark. The
tie-break recovers \textbf{7.9\% of the idle energy on average}
($p<0.001$), with every size group individually significant, and with no
degradation of the primary objective (makespan differences between
$-0.11\%$ and $-0.55\%$ in LexME's favour; we attribute this small,
consistent difference at least partly to machine-load conditions between
the two experimental campaigns and make no claim of makespan improvement
--- the relevant observation is that the tie-break costs nothing).

The recovery grows markedly with instance size and with the job/machine
ratio: from $3.2\%$ on $15{\times}15$ up to $\mathbf{26.5\%}$ on
$50{\times}15$ and $12.6\%$ on $50{\times}20$. The pattern has a clean
explanation: on large instances with a high job/machine ratio the makespan
is comparatively easy to attain (cf.\ the near-zero makespan gaps reported
for $50{\times}15$ in \cite{DiazTabu2026}), so its plateaus are wide, and
wide plateaus are precisely where a blind makespan-only search leaves the
most energy on the table.

\begin{table}[htb]
\centering\small
\caption{LexME vs.\ published makespan-only TS-$N_2$ (per-run means over
30 runs; NPE midpoints; Wilcoxon $p$ paired per instance). RS = NPE after
right-shift timing (heuristic). The spread of the grand mean is best read
from the group rows, which range from $-3.2\%$ to $-26.5\%$; every group is
individually significant.}
\label{tab:lex}
\begin{tabular}{l r rrr rrr rrr}
\toprule
 & & \multicolumn{3}{c}{$\Cmax$ (mid)} & \multicolumn{3}{c}{NPE (mid)} &
 \multicolumn{3}{c}{NPE after RS} \\
\cmidrule(lr){3-5}\cmidrule(lr){6-8}\cmidrule(lr){9-11}
Group & $n$ & base & lex & $\Delta\%$ & base & lex & $\Delta\%$ & base & lex & $\Delta\%$ \\
\midrule
Classical & 12 & 1005.5 & 1003.6 & $-0.19$ & 10978 & 10193 & $\mathbf{-7.1}$ & 9826 & 9228 & $\mathbf{-6.1}$ \\
$15{\times}15$ & 10 & 1254.2 & 1249.9 & $-0.34$ & 26965 & 26102 & $-3.2$ & 24656 & 24128 & $-2.1$ \\
$20{\times}15$ & 10 & 1411.6 & 1407.9 & $-0.26$ & 23330 & 21803 & $-6.5$ & 21105 & 19882 & $-5.8$ \\
$20{\times}20$ & 10 & 1672.3 & 1668.1 & $-0.25$ & 50388 & 48677 & $-3.4$ & 47130 & 45623 & $-3.2$ \\
$30{\times}15$ & 10 & 1851.6 & 1848.0 & $-0.19$ & 18617 & 16690 & $-10.3$ & 16775 & 15128 & $-9.8$ \\
$30{\times}20$ & 10 & 2061.1 & 2054.6 & $-0.31$ & 41629 & 39248 & $-5.7$ & 38437 & 36436 & $-5.2$ \\
$50{\times}15$ & 10 & 2782.6 & 2779.4 & $-0.11$ & 15847 & 11646 & $\mathbf{-26.5}$ & 14252 & 10380 & $\mathbf{-27.2}$ \\
$50{\times}20$ & 10 & 2901.4 & 2885.3 & $-0.55$ & 31214 & 27290 & $-12.6$ & 28852 & 25253 & $-12.5$ \\
\midrule
\emph{Grand mean} & 82 & 1846.5 & 1841.2 & $-0.29$ & 26971 & 24840 & $\mathbf{-7.9}$ & 24756 & 22915 & $\mathbf{-7.4}$ \\
\bottomrule
\end{tabular}

\medskip
\footnotesize All differences are Wilcoxon-significant: $p\le 0.017$ for
$\Cmax$ per group and $p\le 0.008$ for both NPE columns per group;
$p<0.001$ for all grand means.
\end{table}

\subsection{Order vs.\ timing: decomposing the energy gains}
\label{sec:exp-decomp}

The two levers are complementary. Over the 82 instances: the tie-break
(order) contributes $-7.9\%$; right-shifting the baseline's solutions
(timing only) contributes $-8.2\%$; combining both yields
$\mathbf{-15.0\%}$ total idle energy at zero makespan cost (on the
classical group, $-15.9\%$). Crucially, the lexicographic advantage
\emph{survives} timing optimisation: after right-shifting both arms, LexME
remains $-7.4\%$ below the baseline ($p<0.001$; every group individually
significant, up to $-27.2\%$ on $50{\times}15$). The tie-break therefore
finds structurally better orders, not merely orders whose slack timing
could recover anyway.

This decomposition also bounds what timing alone can do: post-hoc
right-shift recovers 7--21\% depending on instance size, far from the
${\sim}56\%$ of the exact fronts at relaxed makespan
(Table~\ref{tab:probe}). The remaining gap is order-mediated at fixed
makespan --- precisely the target of the $N_2^{ME}$ neighbourhood.

\subsection{Validation against exact fronts}
\label{sec:exp-exact}

The $\varepsilon$-constraint probes of Section~\ref{sec:conflict} provide,
on six small and medium instances, exact reference fronts against which
the arms can be positioned. Table~\ref{tab:front} reports the mean gap of
each arm's 30 per-run solutions to the lower envelope of the front,
$\text{gap}\%=(\NPE/\text{front}(\Cmax)-1)\times 100$, where
$\text{front}(\Cmax)$ is the envelope value at the largest probed makespan
allowance not exceeding the solution's makespan. Two caveats apply: the
fronts are exact for the \emph{midpoint scenario} (a crisp proxy of the
interval instance; the mid-point of the interval makespan upper-bounds the
crisp-midpoint makespan by convexity of the longest path, so the lookup is
well defined), and time-limited CP-SAT points are non-monotone, hence the
envelope. Gaps are therefore indicative rather than certificates.

\begin{table}[htb]
\centering\small
\caption{Mean gap to the exact (midpoint-proxy) front, per arm
(30 runs per instance). Lower is better; $0\%$ = on the front. On
\texttt{la29} CP-SAT could not close the tightest allowances and the
envelope starts above the arms' makespan.}
\label{tab:front}
\begin{tabular}{l rrr}
\toprule
Instance & TS-$N_2$ (published) & LexME & LexME+RS \\
\midrule
ft10 & 47.3\% & 39.2\% & 30.2\% \\
la21 & 40.6\% & 23.1\% & 14.6\% \\
la24 & 66.7\% & 48.5\% & 23.5\% \\
la25 & 69.7\% & 58.5\% & 34.5\% \\
tai15\_01 & 25.6\% & 23.6\% & 14.3\% \\
\bottomrule
\end{tabular}
\end{table}

Every rung of the ladder closes part of the distance: the tie-break
removes roughly a quarter of the baseline's excess idle energy and the
right-shift another quarter to a half, leaving the arms within
$14$--$35\%$ of the exact front. That residual gap is order-mediated at
fixed makespan --- precisely the target of the $N_2^{ME}$ neighbourhood
(Section~\ref{sec:n2me}). Figure~\ref{fig:front} illustrates the geometry
on \texttt{la24}: the arms move straight down in the objective space
(constant makespan, decreasing energy) towards the front.

\begin{figure}[htb]
\centering
\begin{tikzpicture}
\begin{axis}[width=0.72\linewidth, height=0.5\linewidth,
  xlabel={$\Cmax$ (midpoint)}, ylabel={NPE (midpoint)},
  legend pos=north east, legend style={font=\footnotesize},
  grid=major, grid style={dotted}]
\addplot[blue, thick, mark=*, mark size=1.6pt]
  table {figures/la24_front.dat};
\addlegendentry{exact front (envelope)}
\addplot[red, only marks, mark=square*, mark size=2.4pt]
  coordinates {(947.8, 6017.9)};
\addlegendentry{TS-$N_2$ (published)}
\addplot[orange, only marks, mark=triangle*, mark size=3pt]
  coordinates {(947.0, 5325.2)};
\addlegendentry{LexME}
\addplot[teal, only marks, mark=diamond*, mark size=3pt]
  coordinates {(947.0, 4422.5)};
\addlegendentry{LexME+RS}
\end{axis}
\end{tikzpicture}
\caption{Exact $\varepsilon$-constraint front (midpoint proxy, lower
envelope) and mean arm positions on \texttt{la24}. Each lever moves the
solution vertically towards the front at essentially constant makespan;
the remaining vertical gap is the target of the energy-aware
neighbourhood. A single lexicographic solution occupies one extreme of
the trade-off; the Pareto method of Section~\ref{sec:pareto} will return
the whole curve.}
\label{fig:front}
\end{figure}

\todo{Optional upgrade (noted as a limitation): recompute the validation
fronts with the exact interval model (shared sequencing booleans, two timing
copies for $p^-/p^+$, LEX2 objectives, as in the multi-component MILP of
\cite{Afsar2022}) rather than the midpoint proxy. The $N_2^{ME}$ outcome on
this instance is reported in Section~\ref{sec:exp-n2me}.}

\subsection{The bi-critical neighbourhood: a scoped result}
\label{sec:exp-n2me}

The LexME-$N_2^{ME}$ arm (identical protocol, energy-aware tabu search)
exhibits a sharp size threshold against LexME-$N_2$. Up to $20{\times}20$
it delivers a further $0.8$--$2.7\%$ NPE reduction (all groups
significant) at a $0.1$--$0.3\%$ makespan cost; from $30{\times}15$
upwards it degrades \emph{both} objectives, up to $+10.3\%$ makespan and
$+91.8\%$ NPE on the $50{\times}$ groups, and triples the arm's total
runtime. Candidate restriction and lazy energy evaluation (each isolated
in dedicated probes) recover part of the loss but do not change the
verdict; the residual mechanism is a search-dynamics one --- energy moves
displace the worsening moves tabu search relies on for diversification,
trapping it on mediocre plateaus. The practical reading is clear and, we
believe, valuable: \emph{the cheap lexicographic tie-break is the robust
winner; specialised energy moves help only where makespan plateaus are
already near-optimal}. Their proper role is re-examined in the
Pareto setting, where covering the front (not protecting the makespan
extreme) is the goal.

\subsection{Front-generation methods: parameter tuning}
\label{sec:exp-tuning}

To compare the three front generators fairly, each is tuned separately
with irace~\cite{Lopez2016}, over the same seven engine parameters used to
tune the tabu search in \cite{DiazTabu2026} (elite size and selection,
max-trials, LS target fraction, tabu bad-iterations, crossover and
mutation operators); the ladder's budget split and cap grid are fixed by
design. Each irace evaluation runs the full method (for the ladder, the
whole anchor-plus-levels pipeline) at the full 900\,s budget and returns the
negated hypervolume of its front against a fixed per-instance reference
point, over the same stratified 20-instance training subset used in
\cite{DiazTabu2026}; each racing run uses a budget of $1\,000$ evaluations
with Friedman-test elimination.

Two design choices deserve justification. First, the tuning subset is
\emph{not} held out from the final evaluation: this is classical parameter
configuration of a randomised heuristic, not supervised learning --- the
tuned object is a handful of engine parameters, not a model of the
instances --- and, decisively, all three arms are tuned on the \emph{same}
subset with the same budget and protocol, so any residual optimistic bias is
shared and cannot alter their comparison; the subset is size-stratified
precisely so that the configuration generalises across the benchmark, and
the group-level results on the remaining 62 instances
(Table~\ref{tab:mo-full}) show no discontinuity between tuning and
non-tuning instances. Second, the per-arm (rather than shared) tuning is
what makes the comparison conservative: each control gets its own best
configuration rather than inheriting ours. The
winning configurations differ markedly across arms
(Table~\ref{tab:tuned-mo}), confirming that per-arm tuning is warranted:
the ladder favours GOX crossover, the more disruptive \emph{swap}
mutation, and lighter per-individual intensification --- consistent with
its need for diversity to spread along the front --- whereas the baselines
retain insertion-based operators.

\begin{table}[htb]
\centering\small
\caption{irace-tuned configurations of the three front generators.}
\label{tab:tuned-mo}
\begin{tabular}{l cccccc}
\toprule
Arm & elite (size/sel) & max-trials & LS target & tabu-iter & crossover & mutation \\
\midrule
LADDER & 40 / 4 & 11 & 0.69 & 16 & GOX & swap \\
ABC-P  & 39 / 2 & 16 & 0.32 & 23 & JOX & insertion \\
MA-P   & 42 / 2 & 24 & 0.86 &  9 & GOX & insertion \\
\bottomrule
\end{tabular}
\end{table}

\subsection{Front quality: decomposition vs.\ dominance across the benchmark}
\label{sec:exp-pareto}

We now put the capped ladder against the two dominance-based controls on the
full 82-instance benchmark (30 runs each, all three arms irace-tuned).
Per instance, we build each arm's global (best-over-runs) front and score it
by hypervolume ratio against a per-instance reference (the non-dominated
union of the three arms plus, where available, the exact front of
Section~\ref{sec:conflict}); statistical testing follows the protocol of
Section~\ref{sec:exp-setup} (Friedman omnibus, Holm-corrected post-hoc,
per-instance Mann--Whitney with Holm).
Table~\ref{tab:mo-full} reports the group-level picture and
Table~\ref{tab:mo} the classical instances in detail.

The headline is that decomposition is the better use of the budget in
aggregate: the ladder attains the highest hypervolume on \textbf{55 of 82}
instances and the best mean ratio overall (\textbf{0.890}, vs.\ $0.749$ for
MA-P and $0.703$ for ABC-P). The Friedman omnibus over per-instance median
per-run HV rejects equality of the three arms
($\chi^2_F=44.8$, $p<10^{-9}$); Holm-corrected post-hoc tests find the
ladder and MA-P each significantly better than ABC-P, and the per-instance
Mann--Whitney counts (Holm-corrected) are $32/82$ ladder-over-ABC-P,
$12/82$ ladder-over-MA-P, with $7$ losses. On the
classical, $15{\times}15$, $20{\times}15$ and $20{\times}20$ groups the
dominance controls are not close (ladder mean $0.875$--$0.965$, winning
$9$--$10$ of every $10$ instances): a standard NSGA-II or Pareto-archive
search leaves a large part of the front unrecovered, which is precisely the
evidence that the front does not come for free and that the decomposition is
doing real work.

The statistics also reveal \emph{where} the ladder's advantage lives, and
it is worth being precise. On the per-run view --- the expected quality of
a \emph{single} 900\,s run --- the ladder and MA-P are globally comparable
(post-hoc $p=0.73$; pooled per-run medians $0.427$ vs.\ $0.410$), both well
above ABC-P ($0.279$). The decisive difference appears in the
\emph{best-over-30-runs} front that a practitioner actually reports: ladder
runs anchor at different makespans and sweep different cap grids, so their
fronts are complementary and their union grows far beyond any single run
($0.890$), whereas MA-P runs converge to similar regions and their union
adds little ($0.749$). Both arms receive identical run budgets and the
union is taken identically, so this is a property of the strategies:
decomposition converts run-to-run variance into coverage; dominance-based
search does not.

The picture is deliberately not a clean sweep, and the exceptions are the
most informative part. On the two largest, machine-heavy groups the ladder's
advantage first vanishes and then reverses: on $30{\times}20$ ABC-P dominates
outright ($0.923$ vs.\ $0.785$, ladder $0/10$), and on the $50{\times}$ groups
MA-P becomes competitive and edges ahead on $50{\times}20$ ($0.854$ vs.\
$0.828$). This boundary has a single, verifiable cause, and it is the same
one that governs the makespan extreme of the ladder. The ladder caps the
whole front at $\kappa=(1+\varepsilon)C^{*}$, where $C^{*}$ is the makespan
of its anchor; a per-run comparison of that anchor against the dedicated
makespan-only tabu search of \cite{DiazTabu2026} shows the two are
statistically equivalent within a $2\%$ margin (two one-sided tests,
$p<10^{-21}$), but with a small median gap of $1.2\%$ that grows with size
and machine count and peaks exactly on the $\times 20$ large groups
($+1.8\%$ on $30{\times}20$, up to $+2.8\%$ on $50{\times}20$). Wherever the
anchor lands above the makespan the free engines can reach, the entire ladder
is confined to higher makespan and cannot populate the low-makespan corner of
the front; the dominance controls, unconstrained, capture that corner and the
hypervolume that comes with it. The diagnosis is direct (Figure~\ref{fig:fronts}):
on \texttt{tai30\_20} instances ABC-P reaches makespans of $2102$/$1992$ against
the ladder anchor's $2137$/$2009$, whereas on \texttt{ft10} the ladder spreads a
deep front the controls never approach. The practical consequence is a sharp
prescription rather than a caveat: \emph{decomposition is the method of choice
whenever a strong makespan anchor is attainable, which is the entire small-
and medium-instance range; on the largest high-ratio instances, either the
anchor must be strengthened} (more anchor budget, or a makespan-specialised
anchor configuration) \emph{or a dominance-based search should be preferred}.
Tuning narrows but does not close the baselines' gap in the ladder's regime
(ABC-P improves from a mean ratio of $0.54$ under the inherited configuration
to $0.70$), so the ordering there is a property of the strategy, not of a
favourable parameterisation.

A word on the shared local search: all three arms use the $N_2^{ME}$
neighbourhood, whose lexicographic-mode pathology on large instances was
documented in Section~\ref{sec:exp-n2me}. Because the component is shared,
it cannot favour one front-generation strategy over another, and in Pareto
mode its energy moves serve front coverage rather than protecting the
makespan extreme; still, an interaction with instance size cannot be
excluded a priori, and a neighbourhood ablation ($N_2$ vs.\ $N_2^{ME}$
inside one arm on the large $\times 20$ groups) is a natural robustness
check for the regime boundary reported here.

\begin{table}[htb]
\centering\small
\caption{Front hypervolume ratio (vs.\ the per-instance union reference;
higher is better) by size group, three tuned arms, 30 runs, best-over-runs
front. ``Wins'' counts instances where the arm attains the highest ratio;
the last block reports per-instance Mann--Whitney significant wins of the
ladder on the 30-run HV distributions, Holm-corrected across the 82
instances ($\alpha=0.05$). Bold: best group mean.}
\label{tab:mo-full}
\begin{tabular}{l r ccc c cc}
\toprule
 & & \multicolumn{3}{c}{mean HV ratio} & wins & \multicolumn{2}{c}{L $>$ (MWU+Holm)} \\
\cmidrule(lr){3-5}\cmidrule(lr){7-8}
Group & \# & LADDER & ABC-P & MA-P & L/A/M & ABC-P & MA-P \\
\midrule
classical      & 12 & \textbf{0.875} & 0.640 & 0.663 & 10/2/0 & 6 & 3 \\
tai $15{\times}15$ & 10 & \textbf{0.960} & 0.750 & 0.756 & 10/0/0 & 7 & 3 \\
tai $20{\times}15$ & 10 & \textbf{0.945} & 0.734 & 0.768 & 9/1/0  & 5 & 3 \\
tai $20{\times}20$ & 10 & \textbf{0.965} & 0.769 & 0.752 & 10/0/0 & 4 & 2 \\
tai $30{\times}15$ & 10 & \textbf{0.888} & 0.721 & 0.705 & 8/1/1  & 2 & 1 \\
tai $30{\times}20$ & 10 & 0.785 & \textbf{0.923} & 0.693 & 0/9/1  & 0 & 0 \\
tai $50{\times}15$ & 10 & \textbf{0.876} & 0.522 & 0.818 & 5/1/4  & 3 & 0 \\
tai $50{\times}20$ & 10 & 0.828 & 0.579 & \textbf{0.854} & 3/2/5 & 5 & 0 \\
\midrule
\textit{All 82} & 82 & \textbf{0.890} & 0.703 & 0.749 & 55/16/11 & 32 & 12 \\
\bottomrule
\end{tabular}
\end{table}

\begin{figure}[htb]
\centering
\begin{tikzpicture}
\begin{axis}[width=0.5\linewidth, height=0.44\linewidth,
  xlabel={$C_{\max}$ (midpoint)}, ylabel={NPE (midpoint)},
  title={(a) \texttt{ft10}, $10{\times}10$},
  legend style={font=\scriptsize, at={(0.98,0.98)}, anchor=north east},
  label style={font=\small}, title style={font=\small},
  tick label style={font=\scriptsize}]
\addplot[blue, thick, mark=*, mark size=1.3pt]
  table[x=cmax, y=npe] {figures/ft10_LADDER.dat};
\addplot[red, mark=square*, mark size=1.7pt]
  table[x=cmax, y=npe] {figures/ft10_ABCP.dat};
\addplot[orange!85!black, mark=triangle*, mark size=2pt]
  table[x=cmax, y=npe] {figures/ft10_MAP.dat};
\legend{LADDER, ABC-P, MA-P}
\end{axis}
\end{tikzpicture}\hfill
\begin{tikzpicture}
\begin{axis}[width=0.5\linewidth, height=0.44\linewidth,
  xlabel={$C_{\max}$ (midpoint)},
  title={(b) \texttt{tai30\_20\_01}, $30{\times}20$},
  label style={font=\small}, title style={font=\small},
  tick label style={font=\scriptsize}, scaled x ticks=false,
  x tick label style={/pgf/number format/1000 sep={}}]
\addplot[blue, only marks, mark=*, mark size=2pt]
  table[x=cmax, y=npe] {figures/tai30_20_01_LADDER.dat};
\addplot[red, only marks, mark=square*, mark size=2.2pt]
  table[x=cmax, y=npe] {figures/tai30_20_01_ABCP.dat};
\addplot[orange!85!black, only marks, mark=triangle*, mark size=2.6pt]
  table[x=cmax, y=npe] {figures/tai30_20_01_MAP.dat};
\end{axis}
\end{tikzpicture}
\caption{Best-over-runs fronts of the three tuned arms (lower-left is better).
\textbf{(a)} On \texttt{ft10} the ladder spreads a deep front that dominates
both controls, reaching low-energy solutions they never approach.
\textbf{(b)} On the large $30{\times}20$ instance the ladder is confined by
its anchor to a narrow high-makespan band ($C_{\max}\!\in[2137,2140]$), while
ABC-P reaches down to $C_{\max}=2102$ and captures the low-makespan corner ---
the mechanism behind the ladder's only systematic regression.}
\label{fig:fronts}
\end{figure}

\begin{table}[htb]
\centering\small
\caption{Per-instance detail of the classical block of
Table~\ref{tab:mo-full}: front hypervolume ratio (vs.\ the per-instance union
reference; higher is better) of the three tuned arms, 30 runs, best-over-runs
front. Bold: best per row.}
\label{tab:mo}
\begin{tabular}{l ccc}
\toprule
Instance & LADDER & ABC-P & MA-P \\
\midrule
ABZ7 & \textbf{0.969} & 0.531 & 0.698 \\
ABZ8 & \textbf{0.947} & 0.607 & 0.537 \\
ABZ9 & 0.936 & \textbf{0.983} & 0.792 \\
FT10 & \textbf{0.816} & 0.667 & 0.760 \\
FT20 & 0.996 & \textbf{0.997} & 0.924 \\
la21 & \textbf{0.904} & 0.687 & 0.888 \\
la24 & \textbf{0.662} & 0.506 & 0.574 \\
la25 & \textbf{0.696} & 0.261 & 0.313 \\
la27 & \textbf{0.993} & 0.565 & 0.760 \\
la29 & \textbf{0.625} & 0.374 & 0.310 \\
la38 & \textbf{0.955} & 0.934 & 0.815 \\
la40 & \textbf{1.000} & 0.564 & 0.586 \\
\midrule
\textit{Mean} & \textbf{0.875} & 0.640 & 0.663 \\
\bottomrule
\end{tabular}
\end{table}

\todo{Optional supplementary material: the additive $\varepsilon$-indicator
against the exact fronts on the probed instances, and the full
tuned-vs-inherited tuning-gain table.}

%=============================================================================
\section{Conclusions}
\label{sec:conclusions}

We introduced the multiobjective green interval JSP, showed that its green
objective reduces to idle passive energy, and that the makespan--energy
conflict only emerges in the (order, timing) decision space --- under
semi-active timing the objectives are almost perfectly aligned, which
explains why a naive biobjective treatment would fail. On this basis we
developed, for a single compromise solution, the lexicographic tie-break,
which harvests the energy hidden in makespan plateaus (7.9\% on average over
82 instances, up to 26.5\% on $50{\times}15$, always at zero makespan cost);
right-shift timing optimisation, which contributes a further ${\sim}8\%$ and
is near-optimally approximated by a single backward pass; and the bi-critical
neighbourhood $N_2^{ME}$, the first sequencing neighbourhood targeting idle
gaps, with feasibility proven and connectivity established for makespan and
machine-completion optimality.

For the Pareto front, our central methodological claim is not the
decomposition idea --- the ladder is the classical $\varepsilon$-constraint
method~\cite{Haimes1971} --- nor that we beat NSGA-II, but that an existing
single-objective interval-JSP solver can be recycled, unchanged, into a strong
front generator by casting each $\varepsilon$-level as a makespan-capped energy
problem and chaining them by warm-starting. With two off-the-shelf
dominance-based algorithms as controls, this capped $\varepsilon$-ladder
attains the highest hypervolume on $55$ of $82$ instances and dominates
throughout the small- and medium-instance range, which is the evidence that a
good front does not come for free on this problem. Just as informative is the
boundary of that advantage: on the largest, machine-heavy instances the
scalar anchor loses a small but decisive amount of makespan quality
(statistically equivalent to the dedicated solver within $2\%$, but enough to
cap the front above the low-makespan corner), and there dominance-based search
becomes competitive again. This yields a mechanism-level prescription for
practice --- decomposition when a strong anchor is attainable, a stronger
anchor or a dominance-based method otherwise --- rather than a single
universal winner. The design is backed throughout by either proof
(goal-programming and $\varepsilon$-constraint equivalence, capped
connectivity) or measurement (the one-move landscape analysis motivating
recombination). Future work follows directly from the boundary we identified:
a makespan-specialised or adaptively-budgeted anchor for large high-ratio
instances, and, more broadly, the flexible interval JSP and vague-goal green
scheduling~\cite{GarciaGomez2026}.

%=============================================================================
\section*{CRediT authorship contribution statement}

\textbf{Hern\'an D\'{\i}az Rodr\'{\i}guez}: Conceptualization, Methodology,
Software, Formal analysis, Investigation, Data curation, Visualization,
Writing -- original draft, Writing -- review and editing.

\section*{Declaration of competing interest}

The author declares that there are no known competing financial interests or
personal relationships that could have appeared to influence the work
reported in this paper.

\section*{Data availability}

\todo{Deposit the run logs, aggregated CSVs, irace records and analysis
scripts on Zenodo and cite the DOI here (mirroring the ASOC manuscript).}

\section*{Acknowledgements}

This work was supported by the Spanish Ministry of Science, Innovation and
Universities (MCIN/AEI/10.13039/501100011033) under grant
PID2022-141746OB-I00.

\bibliographystyle{elsarticle-num}
\bibliography{refs}

\end{document}
